\documentclass[11pt]{article}
\usepackage[margin=1in]{geometry}
\usepackage{amsmath,amssymb,amsthm,mathtools}
\usepackage[hidelinks]{hyperref}
\newtheorem{theorem}{Theorem}
\newtheorem{lemma}[theorem]{Lemma}
\newtheorem{corollary}[theorem]{Corollary}
\title{Uniform adjacent-monomial Strong Factorial by a shared parabola}
\author{Internal algebraic proof; no priority or external-review claim}
\date{17 November 2026}
\begin{document}
\maketitle
\begin{theorem}\label{thm:main}
Let $a\ge0$ be an integer and $\beta>0$. For every nonzero polynomial
$f(X)=X^a(A+BX)$ and every integer $n\ge1$, its two consecutive
Gamma moments at orders $n-1,n$ cannot both vanish. Thus the entire
adjacent-monomial family satisfies Strong Factorial, at every starting
index, for every positive Gamma shape; in particular for ordinary factorials.
\end{theorem}
The theorem does not address arbitrary two-monomial gaps, arbitrary
supports, general Factorial, or unrestricted Strong Factorial.
The prior \href{PADE_ADJACENT_MONOMIALS.tex}{Pad\'e note} proved the ordinary
$a\le6$ ranges and an indexed strip family. Here the uniform conclusion
uses a different separation region. No unbounded vertical spectral bound
for its auxiliary polynomials is asserted or needed.

\section{Two elementary polynomial identities}
For real $t>-1$ define the monic polynomial
\[
 R_N(t,z)=\sum_{k=0}^N(-1)^k\binom Nk(t+1)_k z^{N-k}.
                                                               \tag{1}
\]
Its exponential generating function is
\[
 \sum_{N\ge0}R_N(t,z)\frac{u^N}{N!}
                         =e^{zu}(1+u)^{-t-1}.                  \tag{2}
\]
Coefficient comparison, or differentiation of (2), gives
\[
 (t+1)R_d(t+1,z)=(t+d+1+z)R_d(t,z)-zR_d(t-1,z),              \tag{3}
\]
\[
 R_N(t,z)=zR_{N-1}(t-1,z)-(t+N)R_{N-1}(t,z).                \tag{4}
\]
These identities are polynomial in $t$, so all real parameters in them
are legitimate. For fixed $d,z$, the leading coefficient of $R_d(t,z)$
as a polynomial in $t$ is $(-1)^d$.

\section{Every possible obstruction lies inside a parabola}
Put
\[
 \mathcal P_N=\{z=x+iy:y^2<4N(x+N)\}\qquad(N\ge1).
\]
\begin{lemma}\label{lem:inside}
For every real $M>-1$, every zero of $R_N(M,z)$ belongs to
$\mathcal P_N$.
\end{lemma}
\begin{proof}
Let $S=M+N$ and form the $N$ by $N$ complex symmetric tridiagonal matrix
\[
 T_{jj}=S+2-2j,\qquad
 T_{j,j+1}=T_{j+1,j}=i\sqrt{j(S+1-j)}.                        \tag{5}
\]
The square roots are positive real since $M>-1$. Its leading principal
characteristic polynomials satisfy
\[
 D_0=1,\quad D_1=z-S,\quad
 D_j=(z-S-2+2j)D_{j-1}+(j-1)(S+2-j)D_{j-2}.
\]
The exact coefficient formula
\[
 D_j(z)=\sum_{k=0}^j(-1)^{j-k}\frac{j!}{k!}
                           \binom{S-k}{j-k}z^k
\]
verifies this recurrence, and at $j=N$ gives
$\det(zI-T)=R_N(M,z)$. Generalized binomial coefficients here are
polynomials in their upper argument, so this proof also covers noninteger $M$.

Let $Tu=zu$, normalize $u^*u=1$, and write
$\tau=\sum_{j=1}^N(j-1)|u_j|^2$, so $0\le\tau\le N-1$.
The Hermitian part of $T$ is its real diagonal, hence
\[
 x=S-2\tau,\qquad
 y=2\operatorname{Re}\sum_{j=1}^{N-1}
                  \sqrt{j(S+1-j)}\,\overline{u_j}u_{j+1}.
\]
Cauchy--Schwarz gives
\[
 y^2\le4\Bigl(\sum_{j=1}^{N-1}j|u_{j+1}|^2\Bigr)
           \Bigl(\sum_{j=1}^{N-1}(S+1-j)|u_j|^2\Bigr)
       \le4\tau(S-\tau).
\]
Finally $S>0$, and
\[
 N(x+N)-\tau(S-\tau)
             =(N-\tau)S+(N-\tau)^2>0.                        \tag{6}
\]
This proves the strict parabola inclusion, including $N=1$.
\end{proof}

\section{The complementary zero-free recurrence}
The following elementary principle is the standard continued-fraction
parabola mechanism; compare Saff--Varga, \emph{On the sharpness of
 theorems concerning zero-free regions for certain sequences of polynomials},
Numerische Mathematik 26 (1976), Theorem 1.1, in the
\href{sources/saff-varga-1976-sharpness.pdf}{retained primary source}.
The complete argument needed here follows, so no external theorem or
unreviewed dependency is an input.
\begin{lemma}\label{lem:outside}
Suppose $p_0=1$, $p_{-1}=0$ and
\[
 p_j=(1+z/b_j)p_{j-1}-(z/c_j)p_{j-2},\qquad b_j,c_j>0,
\]
with the second term omitted when $j=1$. Put $b_0=0$.
If $b_j(1-b_{j-1}/c_j)\ge\alpha>0$ for every $j$, then every $p_j$
is nonzero on $\mathcal P_\alpha=\{x+iy:y^2<4\alpha(x+\alpha)\}$.
\end{lemma}
\begin{proof}
Fix $z\in\mathcal P_\alpha$, so $x> -\alpha$.
Inductively put $\mu_j=zp_{j-1}/(b_jp_j)$ and $\mu_0=0$.
If $p_{j-1}\ne0$ and $\operatorname{Re}\mu_{j-1}<1$, write
\[
 d_j=b_{j-1}b_j/c_j,\quad v=b_j-d_j\mu_{j-1},\quad
                  \mu_j=\frac{z}{z+v}.
\]
Here $\operatorname{Re}v\ge b_j-d_j\ge\alpha$, and
$\operatorname{Re}(z+v)>0$, so the next denominator is nonzero.
For $v=u+iw$, $u\ge\alpha$, we have
\[
 |v|^2+\operatorname{Re}(v\overline z)
 =u^2+ux+(w+y/2)^2-y^2/4
 \ge\alpha^2+\alpha x-y^2/4>0.                              \tag{7}
\]
The inequality uses $2u+x>0$ for $u\ge\alpha$.
It implies $\operatorname{Re}\mu_j<1$ because
$1-\operatorname{Re}\mu_j=operatorname{Re}(v/(z+v))$.
The induction starts at $p_0=1,\mu_0=0$ and proves all assertions.
\end{proof}

\begin{theorem}\label{thm:coprime}
Let $N\ge2$, $l>-1$ be real, and let $M=l+a$ with $a\ge1$ an integer.
Then $R_N(M,z)$ and $R_{N-1}(l,z)$ have no common complex zero.
\end{theorem}
\begin{proof}
Their constant terms are nonzero by (1), so a common zero $z$ is nonzero.
Put $S_t=R_{N-1}(t,z)$. Equations (3)--(4) imply
\[
 S_l=0,\qquad zS_{M-1}=(M+N)S_M.                            \tag{8}
\]
If $a=1$, these relations force $S_l=S_{l+1}=0$; recurrence (3)
then propagates zeros to every larger parameter $t=l+k$.
That contradicts the nonzero leading coefficient of $S_t$ as a polynomial
in $t$. The same reasoning shows that, for $a\ge2$, the vector
$(S_{l+1},\ldots,S_{M-1})$ is nonzero.

Set $D=a-1$. Recurrence (3), with the two boundary conditions (8),
shows that this nonzero vector lies in the kernel of a tridiagonal
pencil. Its diagonal is
\[
 l+N+j+z\quad(1\le j<D),\qquad
                   l+N+D+\frac{N}{M+N}z\quad(j=D),
\]
its upper entry in row $j$ is $-(l+j+1)$, and its lower entry is $-z$.
The last diagonal follows by replacing $S_M$ with
$zS_{M-1}/(M+N)$ in row $M-1$ of (3).
Its determinant therefore vanishes.

Normalize each leading principal determinant to constant term one,
calling the result $p_j$. For $j<D$ the coefficients in
Lemma~\ref{lem:outside} are
\[
 b_j=l+N+j,\qquad
 c_j=\frac{(l+N+j)(l+N+j-1)}{l+j}\quad(j\ge2).
\]
For the last row they are
\[
 b_D=\frac{(M+N)(M+N-1)}{N},\qquad
 c_D=\frac{(M+N-1)(M+N-2)}{M-1}\quad(D\ge2).
\]
All quantities are positive: $l>-1$, and $M>1$ when $a\ge2$.
When $D=1$ only $p_1=1+z/b_1$ occurs, with the displayed last-row $b_1>N$.
When $D\ge2$, the first value is $b_1>N$, the interior values
$b_j(1-b_{j-1}/c_j)$ equal $N$, and the last value is $M+N>N$.
Thus Lemma~\ref{lem:outside}, with $\alpha=N$, forbids a determinant zero
in $\mathcal P_N$. Lemma~\ref{lem:inside} places our common zero there.
This contradiction proves the theorem for every $a$, with no cutoff.
\end{proof}

\section{Transfer to factorial moments and the exact scope}
For Gamma shape $\beta>0$, the normalized $n$th moment of
$X^a(A+BX)$ is, when $A\ne0$,
\[
 \mathcal L_\beta(f^n)=A^n(\beta)_{an}B_{n,a,\beta}(B/A),\qquad
 B_{n,a,\beta}(t)=\sum_{k=0}^n\binom nk(\beta+an)_k t^k.
\]
Exactly
\[
 R_n(\beta+an-1,z)=z^n B_{n,a,\beta}(-1/z).                  \tag{9}
\]
For $a\ge1,n\ge2$, apply Theorem~\ref{thm:coprime} with
$N=n$, $M=\beta+an-1$, $l=\beta+a(n-1)-1>-1$.
A nonzero ratio $B/A$ produces the same nonzero $z=-A/B$ at both
orders, so (9) excludes simultaneous vanishing. If a coefficient is
zero, the polynomial is a monomial with every moment nonzero.
At terminal $n=1$ the order-zero moment is one.

For $a=0$, let $T_n(z)=R_n(\beta-1,z)$. Formula (2) gives
\[
 T_{n+1}=(z-\beta-n)T_n+znT_{n-1},\qquad T_0=1.
\]
A common zero of consecutive members is nonzero, since
$T_n(0)=(-1)^n(\beta)_n\ne0$. The recurrence propagates that common
zero backward to $T_0$, a contradiction. This completes
Theorem~\ref{thm:main}.

A direct multivariate corollary is the same two-moment assertion for
$\prod_j X_j^{a_j}(A+BX_1)$ under independent positive Gamma shapes:
all other variables supply nonzero rising-factorial factors at each
moment. This corollary requires the monomials to differ by exactly one
power in a single variable; it makes no assertion for general gaps or
multivariate supports.

\section{Exact recurrence at the next, nonadjacent boundary}
For a positive integer gap $d$ define
\[
 U_N^{(d)}(t,z)=\sum_{k=0}^N(-1)^k\binom Nk(t+1)_{dk}z^{N-k}.
\]
For $f=X^a(A+BX^d)$, its $N$th moment has the same reciprocal
normalization with $U_N^{(d)}(\beta+aN-1,z)$.
Direct coefficient comparison gives
\[
 U_N^{(d)}(t,z)=zU_{N-1}^{(d)}(t,z)
                    -(t+1)_d U_{N-1}^{(d)}(t+d,z),           \tag{10}
\]
\[
 (t+1)U_N^{(d)}(t+1,z)
 =(t+1+dN)U_N^{(d)}(t,z)-dzN U_{N-1}^{(d)}(t,z).            \tag{11}
\]
Let $V_t=U_{N-1}^{(d)}(t,z)$. Eliminating $U_{N-2}^{(d)}$
from (10)--(11), for $N\ge2$, yields the sparse order-$d+1$ recurrence
\[
 z(V_t-V_{t+1})=(t+2)_{d-1}
       \bigl((t+dN+1)V_{t+d}-(t+d+1)V_{t+d+1}\bigr).         \tag{12}
\]
The $N=1$ identity is immediate as well.
A hypothetical adjacent obstruction, with
$l=\beta+a(N-1)-1$, $M=l+a$, supplies
\[
                    V_l=0,\qquad zV_M=(M+1)_dV_{M+d}.       \tag{13}
\]
For $d=1$, (12) is exactly the second-order parameter recurrence
used above, and (13) closes its finite continuant problem.
For $d\ge2$, recurrence (12) has order $d+1$ and these are still only
two boundary conditions. They do not by themselves define a square
finite-matrix eigenproblem. Additional structure of this particular
polynomial solution is needed before applying a zero-free-region
argument. Equations (10)--(13) are exact reusable identities, not a
coprimality, squarefreeness or SFC theorem for a nonadjacent gap.

The universal evidence is the elementary inequalities and polynomial
identities above. The accompanying checker replays identities, the
pencil normalization and representative exact rational parameters;
finite replay alone is not the theorem. Independent replay, formal
verification, external review, priority claims, general FC and unrestricted
SFC closure all remain false. The former uniform vertical-bound target
may remain unproved without leaving a gap in this parabola proof.
\end{document}
